\documentclass[11pt]{article}
\usepackage[margin=1in]{geometry}
\usepackage{amsmath,amssymb,amsthm}
\usepackage[T1]{fontenc}
\usepackage[utf8]{inputenc}
\usepackage{hyperref}
\hypersetup{colorlinks=true,linkcolor=blue,citecolor=blue,urlcolor=blue}
\usepackage{booktabs}
\usepackage{enumitem}
\usepackage{array}

\newtheorem{theorem}{Twierdzenie}
\newtheorem{axiom}{Aksjomat}
\newtheorem*{remark}{Uwaga}
\newtheorem{proposition}{Wniosek}

\title{\textbf{$\Sigma$: hierarchiczne, funkcjonalnie ważone\\ rozszerzenie zintegrowanej teorii informacji}}
\author{[Autor]\\ \small w dialogu z Claude (Anthropic)}
\date{sierpień 2026 --- wersja robocza}

\begin{document}
\maketitle

\begin{abstract}
Zintegrowana teoria informacji (IIT) sprowadza stopień świadomości systemu do jednej liczby, $\Phi$, wyznaczanej przez minimalizację miary odległości po partycjach systemu. To zaciera dwa osobne pytania, na które IIT nie odpowiada z osobna: jak bardzo zintegrowany jest system, oraz gdzie w jego architekturze ta integracja ma funkcjonalne znaczenie dla samo-modelowania. Proponujemy $\Sigma$ --- hierarchiczne rozszerzenie, w którym $\Phi$ liczone jest jako profil po warstwach funkcjonalnych $\ell=1,\dots,L$ i łączone przez wagę funkcjonalną $w(\ell)$, inspirowaną postulatem Ag\"uera y Arcasa, że o statusie systemu decyduje to, \emph{co} on oblicza, nie tylko jego struktura przyczynowo-skutkowa. Wprowadzamy ponadto dynamiczny człon napędzany błędem samo-predykcji, odróżniający samo-modelowanie deskryptywne od konstytutywnego (to drugie wymaga aktualizacji parametrów online, nie tylko płynnej autonarracji), i formalizujemy stan systemu jako trójkę $(\Sigma,\Gamma,A)$: zintegrowana informacja, aktywność konstytutywnej samo-korekty, oraz ich przestrzenne wyrównanie. Dowodzimy własności brzegowych wszystkich trzech wielkości, rozszerzamy model o miękki mechanizm konkurencji/wykluczenia dla nakładających się podsystemów (ze świadomym, uzasadnionym odejściem od kanonicznego, binarnego postulatu wykluczenia IIT), przeliczamy przykład numeryczny i odnosimy formalizm do współczesnych architektur AI, ostrząc rozróżnienie między konstytutywnym modelowaniem świata (jak w architekturach pamięci testowo-czasowej) a właściwym konstytutywnym samo-modelowaniem. To jest robocza synteza koncepcyjna, niezweryfikowana recenzyjnie; dołączony jest wstępny przegląd literatury.
\end{abstract}

\section{Wstęp}

Standardowe IIT definiuje zintegrowaną informację przez
\[
\varphi_P(s) = D\big(p(s'\mid s)\,\|\,p(s'\mid s)|_P\big), \qquad \Phi(s) = \min_{P\in\mathcal P}\varphi_P(s),
\]
gdzie $D$ to miara odległości (np. earth mover's distance) między faktycznym rozkładem przejść systemu a rozkładem po przecięciu go wg partycji $P$. Powracająca krytyka (bliska nieformalnemu zarzutowi Aaronsona \cite{aaronson2014} i bardziej formalnemu zarzutowi Cerullo \cite{cerullo2015}) mówi, że $\Phi$ zaciera bogatą strukturę przyczynowo-skutkową do jednej liczby, bez rozróżnienia \emph{rodzaju} czy \emph{lokalizacji} integracji. Sama miara bazowa jest sporna: $\Phi$ nie jest dobrze zdefiniowane dla ogólnych układów fizycznych \cite{barrett2019}, a warianty $\Phi$ silnie się rozjeżdżają w symulacji \cite{mediano2019}. $\Sigma$ dziedziczy tę zależność. Jest to koncepcyjnie pokrewne, lecz odrębne od nieredukowalności obliczeniowej Wolframa --- idei, że niektóre obliczenia nie dają się skrócić i trzeba je przeprowadzić krok po kroku, by poznać wynik.

Ta praca rozwija hierarchiczne, funkcjonalnie ważone rozszerzenie, $\Sigma$, motywowane trzema dodatkowymi względami: (i) różne warstwy funkcjonalne systemu (homeostaza, percepcja/motoryka, model świata, model siebie) prawdopodobnie różnie przyczyniają się do podmiotowości nawet przy równym $\Phi$; (ii) płynność autonarracji nie jest dowodem realnej pętli samo-modelującej (literatura o konfabulacji przy rozdzielonych półkulach, via interpreter Gazzanigi, to paradygmatyczny przypadek); (iii) konkurujące, nakładające się kandydackie podsystemy mogą nie tylko \emph{nie istnieć} (jak w postulacie wykluczenia IIT), ale aktywnie degradować integrację systemu nadrzędnego --- możliwość, której binarne wykluczenie IIT nie potrafi wyrazić.

\subsection*{Wkład pracy}
\begin{itemize}[leftmargin=*]
\item Hierarchiczny profil $\Phi(s,\ell,t)$ łączony przez wagę funkcjonalną $w(\ell)$ w $\Sigma(t)$, z dokładnym rozkładem jego pochodnej czasowej.
\item Człon błędu samo-predykcji $e_\ell(t)$ i bramka konstytutywności $g_\ell(t)$ formalizujące rozróżnienie deskryptywne/konstytutywne, dające drugą wielkość $\Gamma(t)$ niezależną od $\Sigma$ z konstrukcji.
\item Trzecia diagnostyka, $A(t)\in[-1,1]$, mierząca przestrzenne wyrównanie między integracją a samo-korektą, z dowodem ograniczenia przez nierówność Cauchy'ego--Schwarza.
\item Dowiedzione własności brzegowe wszystkich trzech wielkości, w tym jawny aksjomat ($w_0\ge 0$) potrzebny dla $\Sigma\ge 0$.
\item Miękki mechanizm konkurencji/wykluczenia dla nakładających się, przeszkadzających sobie warstw, z dokładnym progiem krytycznym, i jawne, uzasadnione odejście od kanonicznego, binarnego wykluczenia IIT.
\item Przeliczony przykład numeryczny oraz odniesienie formalizmu do współczesnych architektur AI.
\end{itemize}

\section{Prace pokrewne}

\textbf{IIT 4.0} \cite{albantakis2023} formalizuje wykluczenie jako ściśle binarne: spośród nakładających się kandydackich substratów istnieje tylko ten o maksymalnym $\varphi_s$; pozostałe są całkowicie wykluczone (nie wnoszą nic), nie stopniowo tłumione. To istotne dla \S\ref{sec:exclusion}.

\textbf{Nested Observer Windows} \cite{riddle2024} niezależnie proponują liczenie $\Phi$ osobno na zagnieżdżonych poziomach hierarchii i badanie stosunku między poziomami jako testu między odczytem panpsychistycznym a emergentystycznym --- bliskie w duchu naszemu warstwowemu profilowi $\Phi(s,\ell,t)$, choć bez ważenia funkcjonalnego i dynamiki.

\textbf{Consciousness as Uncommon Self-Knowledge} \cite{usk2026} proponuje, że liczy się tylko \emph{synergiczny, samo-skierowany} komponent informacji, przewidując, że systemy o wysokim $\Phi$ ale bez self-modelu wypadają blisko zera pod ich miarą --- niezależna zbieżność z intuicją stojącą za naszym rozdziałem $\Gamma/A$.

\textbf{Predictive processing i active inference} \cite{friston2010,deane2021} motywują formalizm błędu samo-predykcji; adwersarialna współpraca IIT vs. predictive processing \cite{intrepid2026} potwierdza, że połączenie obu pozostaje otwartym pytaniem badawczym.

\textbf{Teorie wyższego rzędu} \cite{rosenthal2005} służą jako punkt kontrastowy: nasze rozróżnienie deskryptywne/konstytutywne adresuje podobne pytanie (czy autonarracja stanu mentalnego jest dowodem realnego dostępu wyższego rzędu?) przez kryterium funkcjonalne (czy błąd napędza aktualizację?), nie przez samo istnienie reprezentacji wyższego rzędu.

Postulat Ag\"uera y Arcasa, że żywa materia jest materią \emph{funkcjonalną} --- liczy się to, co system oblicza, nie tylko jego struktura przyczynowo-skutkowa --- bezpośrednio motywuje wagę funkcjonalną $w(\ell)$.

\textbf{Wielowymiarowa świadomość.} Bayne, Hohwy \& Owen \cite{bayne2016} argumentują, że świadomość lepiej opisać kilkoma wymiarami niż jednym ,,poziomem''; $(\Sigma,\Gamma,A)$ to ugruntowana w IIT, dynamiczna instancja tego programu. \textbf{Integrated World Modeling Theory} \cite{safron2020} łączy IIT, Global Neuronal Workspace i zasadę wolnej energii --- najbliższy odpowiednik ambicji $\Sigma$, ale jako architektura koncepcyjna, nie formalizm z dowiedzionymi własnościami brzegowymi. \textbf{Postulat wykluczenia} jest niezależnie problematyczny: czyni $\Phi$ nie-jednoznacznym \cite{hanson2023}, co jest dodatkowym powodem, by go zastąpić, a nie reinterpretować (\S\ref{sec:exclusion}).

Nie znamy wcześniejszej pracy łączącej hierarchiczne $\Phi$, ważenie funkcjonalne, dynamikę napędzaną błędem samo-predykcji i miękkie wykluczenie z dokładnym progiem w jeden formalizm z dowiedzionymi własnościami brzegowymi; ta synteza wydaje się oryginalna, choć każdy komponent ma niezależny precedens.

\section{Formalizm $\Sigma$}

\subsection{Profil hierarchiczny i waga funkcjonalna}
System $S$ dzielimy na warstwy funkcjonalne $\ell=1,\dots,L$ (np. homeostaza $\to$ percepcja/motoryka $\to$ model świata $\to$ model siebie). Indeks warstwy jest z natury dyskretny: $L$ to mała, skończona liczba całkowita --- kilka pięter funkcjonalnych, nie continuum --- i nie ma granicy $L\to\infty$. Zamiast jednej liczby otrzymujemy profil $\Phi(s)=(\Phi_1(s),\dots,\Phi_L(s))$, zwinięty przez wagę funkcjonalną w skończoną sumę
\[
\Sigma(t) = \sum_{\ell=1}^{L} w(\ell,t)\,\Phi(s,\ell,t).
\]
Każde $\Phi_\ell$ w praktyce szacowano by przez wykonalny proxy, np. miarę dekodowalności z \cite{oizumi2016}. Człony międzywarstwowe $\Phi_{\ell,\ell-1}$ (sprzężenie strukturalne między sąsiednimi warstwami, liczone tą samą miarą odległości co bazowe $\Phi$) są nieujemne z konstrukcji, będąc same odległością.

\subsection{Dynamika czasowa}
Różniczkując $\Sigma(t)$ przy ogólnej (czasowo zmiennej) wadze otrzymujemy
\[
\frac{d\Sigma}{dt} = \sum_{\ell=1}^{L} \left[\frac{\partial w(\ell,t)}{\partial t}\,\Phi(s,\ell,t) + w(\ell,t)\,\frac{\partial \Phi(s,\ell,t)}{\partial t}\right],
\]
rozdzielając zmianę wynikającą ze zmiany istotności funkcjonalnej od zmiany samej integracji.

\subsection{Błąd samo-predykcji i bramka konstytutywności}
Zgodnie z rozróżnieniem deskryptywne/konstytutywne (płynna autonarracja, jak w interpreterze Gazzanigi przy rozdzielonych półkulach, nie jest dowodem realnej pętli zwrotnej), definiujemy błąd samo-predykcji
\[
e_\ell(t) = \big\|\hat s_\ell(t+\Delta t\mid t) - s_\ell(t+\Delta t)\big\|^2 \ge 0
\]
(zero dla warstw bez mechanizmu samo-predykcji), oraz znormalizowaną bramkę konstytutywności
\[
g_\ell(t) = \left|\frac{\partial \theta_\ell(t)}{\partial e_\ell(t)}\right| \in [0,1],
\]
gdzie $g_\ell=0$ oznacza warstwę liczącą błąd, ale nic z nim nie robiącą (deskryptywne), a $g_\ell=1$ --- pełne sprzężenie konstytutywne.

\subsection{Stan jako trójka $(\Sigma,\Gamma,A)$}
Waga jest statyczna, $w(\ell,t)=w_0(\ell)$. Stan systemu to wtedy trójka skalarów zbudowanych z rozłącznych zmiennych:
\[
\Sigma(t) = \sum_{\ell=1}^{L} w_0(\ell)\,\Phi(s,\ell,t), \qquad \Gamma(t) = \sum_{\ell=1}^{L} g_\ell(t)\,e_\ell(t).
\]
$\Sigma$ to architektonicznie ważona zintegrowana informacja; $\Gamma$ to całkowita aktywność konstytutywnej samo-korekty. Ponieważ $w_0$ nie zależy od $g_\ell$ ani $e_\ell$, obie nie dzielą żadnego składnika: każda zmierzona korelacja między nimi w realnym systemie jest odkryciem empirycznym, nie artefaktem definicyjnym. (Waga sprzężona jako $w_0(\ell)+\beta g_\ell(t)e_\ell(t)$ sprawiłaby, że $\Sigma$ i $\Gamma$ korelują tautologicznie.) Świadomość jako kilka wymiarów, a nie jedna liczba, idzie za \cite{bayne2016}; $(\Sigma,\Gamma,A)$ to ugruntowana w IIT, dynamiczna instancja tego programu.

Trzecia wielkość wychwytuje, czy samo-korekta zachodzi \emph{tam}, gdzie jest integracja:
\[
A(t) = \frac{\mathrm{Cov}_\ell(t)}{\sigma_\Phi(t)\sigma_{ge}(t)} \in [-1,1], \qquad \mathrm{Cov}_\ell(t)=\frac{1}{L}\sum_{\ell=1}^{L}\big[\Phi(s,\ell,t)-\bar\Phi(t)\big]\big[g_\ell(t)e_\ell(t)-\overline{ge}(t)\big],
\]
z konwencją $A(t):=0$, gdy któryś z profili jest jednorodny (wariancja zero). Ujemne $A$ mimo wysokich $\Sigma,\Gamma$ to architektoniczna sygnatura interpretera Gazzanigi: płynna, aktywna samo-korekta zamknięta w płytkiej warstwie, odcięta od integracyjnego rdzenia systemu.

\section{Dowiedzione własności brzegowe}

\begin{axiom}[A1]
$w_0(\ell)\ge 0$ dla wszystkich $\ell$.
\end{axiom}

\begin{theorem}[Nieujemność $\Sigma$]
Przy A1, $\Sigma(t)\ge 0$.
\end{theorem}
\begin{proof}
$\Phi(s,\ell,t)\ge 0$ (odległość) i $w_0(\ell)\ge 0$ z A1; suma iloczynów składników nieujemnych jest nieujemna.
\end{proof}

\begin{theorem}[Nieujemność $\Gamma$, bezwarunkowa]
$\Gamma(t)\ge 0$ zawsze.
\end{theorem}
\begin{proof}
$g_\ell(t)\in[0,1]$ i $e_\ell(t)\ge 0$ (norma do kwadratu); żaden aksjomat nie jest potrzebny.
\end{proof}

\begin{theorem}[Ograniczenie $A$]
$A(t)\in[-1,1]$, z równością wtedy i tylko wtedy, gdy $g_\ell e_\ell$ jest rosnącą (odp. malejącą) funkcją afiniczną $\Phi(s,\ell,t)$ po $\ell$.
\end{theorem}
\begin{proof}
Niech $f(\ell)=\Phi(s,\ell,t)-\bar\Phi(t)$ oraz $h(\ell)=g_\ell(t)e_\ell(t)-\overline{ge}(t)$. Wtedy $\frac1L\sum_{\ell}(\sigma_{ge}f(\ell)-\sigma_\Phi h(\ell))^2 \ge 0$ rozwija się do $2\sigma_{ge}^2\sigma_\Phi^2 \ge 2\sigma_\Phi\sigma_{ge}\,\mathrm{Cov}_\ell$, dając $\mathrm{Cov}_\ell\le\sigma_\Phi\sigma_{ge}$, czyli $A\le1$; symetryczny argument z $\frac1L\sum_{\ell}(\sigma_{ge}f(\ell)+\sigma_\Phi h(\ell))^2\ge0$ daje $A\ge-1$. Równość zachodzi dokładnie wtedy, gdy oba profile są powiązane afinicznie (warunek równości Cauchy'ego--Schwarza).
\end{proof}

\begin{theorem}[Monotoniczność]
$\partial\Gamma/\partial g_\ell = e_\ell\ge0$, $\partial\Gamma/\partial e_\ell=g_\ell\ge0$, oraz $\partial\Sigma/\partial\Phi(\ell)=w_0(\ell)\ge0$ przy A1: brak kontrintuicyjnych efektów znaku w reakcji którejkolwiek wielkości na swoje wejścia.
\end{theorem}

\begin{remark}[Dyskretność]
$\Sigma$, $\Gamma$ i $A$ to skończone sumy po ustalonym, niewielkim zbiorze warstw funkcjonalnych. Oś warstw nie jest continuum: nie ma idealizacji $L\to\infty$ ani postaci całkowej. (Wcześniejsza wersja robocza ją miała; jej zachowanie przy $L\to0$, $\Sigma\to0$ niezależnie od tego, jak skoncentrowana jest integracja, było czystym artefaktem znikającej miary całkowania i tutaj nie występuje.)
\end{remark}

\subsection{Dynamika dyskretna: asymetria między $\Sigma$ a $\Gamma$}
Ponieważ $w_0(\ell)$ jest statyczne w wersji kanonicznej, $\Delta\Sigma=\sum_\ell w_0(\ell)\Delta\Phi(\ell)$ dokładnie, bez żadnego członu korekcyjnego, dla dowolnie dużych (nieinfinitezymalnych) kroków czasowych. $\Gamma=\sum_\ell g_\ell e_\ell$ jest jednak iloczynem dwóch wielkości czasowych, więc
\[
\Delta(g_\ell e_\ell) = \Delta g_\ell\cdot e_\ell(t_0) + g_\ell(t_0)\cdot\Delta e_\ell + \Delta g_\ell\cdot\Delta e_\ell,
\]
a człon krzyżowy znika tylko w granicy infinitezymalnej. W symulowanym scenariuszu dysocjacji (\S\ref{sec:example}) ten człon odpowiadał za około jedną trzecią całkowitej zmiany na pojedynczej warstwie --- niepomijalne dla symulacji w skali zdarzeń dyskretnych (np. epizodu dysocjacyjnego), gdzie ciągłe przybliżenie $\partial\Gamma/\partial t$ systematycznie nie doszacowałoby zmienności $\Gamma$.

\section{Warstwy konkurujące: miękki mechanizm wykluczenia}\label{sec:exclusion}

Czy warstwa może aktywnie \emph{degradować} globalną integrację, zamiast po prostu nie wnosić wkładu (metaforyczna paralela do niejednoznaczności znaku stałej kosmologicznej)? Kandydaci: konkurujące podnarracje przy dysocjacji, lub interpreter Gazzanigi działający w trybie \emph{nadpisującym}, zniekształcającym, nie tylko biernie-deskryptywnym.

Rozbijamy wagę na dwie nieujemne składowe, $w(\ell)=w_+(\ell)-w_-(\ell)$, uziemione w postulacie wykluczenia IIT (wygrywa maksymalnie nieredukowalny kandydat):
\[
\ell^*(t) = \arg\max_\ell \Phi(s,\ell,t), \qquad \text{overlap}(\ell,\ell^*) = \frac{|S_\ell\cap S_{\ell^*}|}{|S_\ell\cup S_{\ell^*}|}\in[0,1]\ \text{(indeks Jaccarda na substratach)},
\]
\[
w_-(\ell,t) = \begin{cases} 0 & \ell=\ell^*(t) \\ \gamma\cdot\text{overlap}(\ell,\ell^*(t))\cdot\Phi(s,\ell,t) & \ell\neq\ell^*(t)\end{cases}, \quad \gamma\ge0.
\]

\begin{proposition}[Dokładny próg liniowy]
$\Sigma(\gamma) = \Sigma_+ - \gamma K$, gdzie $K=\sum_{\ell\neq\ell^*}\text{overlap}(\ell,\ell^*)\Phi(\ell)^2\ge0$ (uwaga na kwadrat: wkład odejmowany od warstwy to $w_-(\ell)\Phi(\ell)=\gamma\cdot\text{overlap}\cdot\Phi(\ell)^2$), a $\Sigma_+$ to wartość bezkonkurencyjna ($\gamma=0$). Netto fragmentacja ($\Sigma<0$) jest osiągalna \emph{wtedy i tylko wtedy, gdy} $\gamma>\gamma_{\text{crit}}=\Sigma_+/K$ --- jedno porównanie skalarne, nie widełki (jeśli $K=0$, żadna warstwa niebędąca zwycięzcą nie nakłada się na $\ell^*$, więc $\Sigma\equiv\Sigma_+$ i fragmentacja jest nieosiągalna). Zastąpienie każdego overlap jego maksimum $1$ daje tylko luźne widełki $\Sigma_+-\gamma\sum_{\ell\neq\ell^*}\Phi(\ell)^2\le\Sigma(\gamma)\le\Sigma_+$, ciasne jedynie dla w pełni zagnieżdżonych substratów.
\end{proposition}

\begin{remark}[Nieciągłość przy przełączeniu $\ell^*$]
$\ell^*(t)$ to $\arg\max$, więc gdy zmienia się warstwa maksymalna, $\ell^*$ --- a z nią $w_-$ i $\Sigma(t)$ --- skacze (\S\ref{sec:example} pokazuje takie przełączenie). To jest zamierzone: nagła zmiana tego, który kompleks ,,wygrywa'', to dokładnie zdarzenie dysocjacyjne, które mechanizm ma reprezentować, a nie nieciągłość do wygładzenia.
\end{remark}

\subsection*{Odrzucona, wierniejsza IIT alternatywa}
Sformalizowaliśmy też wersję uziemioną bezpośrednio w rzeczywistym (binarnym) mechanizmie wykluczenia IIT: binarna bramka nakładania się plus gładka relaksacja samej konkurencji, $s(\ell,t)=\tanh\!\big(\tfrac{\kappa}{2}[\Phi(\ell^*,t)-\Phi(\ell,t)]\big)$, tłumiąca bezpośrednio $\Phi$ (jak w kanonicznym IIT), nie wagę. Ten mechanizm dowodliwie wyklucza $\Sigma<0$ (czynnik tłumienia leży w $(0,1]$) i dokładnie odtwarza kanoniczne, twarde wykluczenie przy $\kappa\to\infty$. Świadomie go odrzucamy: zdolność modelu do wyrażenia netto fragmentacji uznajemy za cenniejszą niż ścisła wierność binarnemu postulatowi IIT, dlatego mechanizm $w_+-w_-$ powyżej pozostaje kanoniczny. To świadome, uzasadnione odejście od IIT, nie jego nieporozumienie, i tak należy je nazwać w każdej przyszłej pracy.

\begin{remark}[Co twierdzi $\Sigma<0$]
Miary zintegrowanej informacji typu ,,całość minus części'' już przyjmują wartości ujemne \cite{barrett2019,phiid2019}, gdzie ujemne oznacza, że całość jest \emph{redukowalna} (netto redundancja). $\Sigma<0$ tutaj to co innego: żadne $\Phi_\ell$ nie schodzi poniżej zera (każde $\Phi_\ell\ge0$); deficyt bierze się z tego, że anti-waga $w_-$ konkurującej warstwy odejmuje więcej wagi architektonicznej niż wnoszą pozostałe warstwy. Deklarowaną nowością jest mechanizm anti-wagi i jego dokładny próg, nie ujemna zintegrowana informacja jako taka. To, że sam postulat wykluczenia jest formalnie problematyczny \cite{hanson2023}, jest niezależnym wsparciem dla jego zastąpienia.
\end{remark}

\section{Przykład numeryczny}\label{sec:example}

Trzy warstwy (homeostaza, model świata, model siebie) z $\Phi=(0{,}2;0{,}6;0{,}9)$, $w_0=(0{,}1;0{,}3;0{,}6)$, $e=(0;0{,}4;0{,}8)$, $g=(0;0{,}3;0{,}9)$.

\begin{center}
\begin{tabular}{@{}lccc@{}}
\toprule
& $\Sigma$ & $\Gamma$ & $A$ \\
\midrule
Bez konkurencji & 0,740 & 0,840 & 0,90 \\
Z konkurencją ($\gamma=0{,}5$, overlap$(2,3)=0{,}4$) & 0,668 & --- & --- \\
\bottomrule
\end{tabular}
\end{center}
Przy overlapach Jaccarda $\{0;0{,}4\}$, $K=0{,}4\times0{,}6^2=0{,}144$, więc dokładna tożsamość daje $\Sigma(0{,}5)=0{,}740-0{,}5\times0{,}144=0{,}668$ (luźne widełki $[0{,}34;0{,}74]$ też ją zawierają, ale są niedopięte, bo $\text{overlap}(1,3)=0$), a $\gamma_{\text{crit}}=0{,}740/0{,}144\approx5{,}14$ --- potwierdzone numerycznie ($\Sigma(5{,}14)\approx0$, $\Sigma(6)=-0{,}124$). Netto fragmentacja wymaga tu siły konkurencji $\sim5\times$ bazowej wagi modelu siebie.

Dwuchwilowy scenariusz dysocjacji (konkurująca podnarracja w warstwie 2 przejmuje rolę $\ell^*$ od warstwy 3) daje $(\Sigma,\Gamma,A)$: $(0{,}740;0{,}840;0{,}90)\to(0{,}575;0{,}680;0{,}96)$. Zauważalnie $A$ \emph{rośnie}, mimo że $\Sigma,\Gamma$ obie maleją --- ilustrując, że $A$ odpowiada na inne pytanie (czy samo-korekta nadąża za aktualnym centrum integracyjnym?) niż pytania o wielkość, na które odpowiadają $\Sigma,\Gamma$.

\section{Zastosowanie do architektur transformerowych}

Standardowy, zamrożony transformer przy inferencji ma $g_\ell(t)\equiv0$ (brak aktualizacji parametrów online), stąd $\Gamma(t)\equiv0$ z definicji, niezależnie od płynności autonarracji, a $A(t):=0$ z konwencji brzegowej (brak wariancji w tożsamościowo zerowym profilu). Dostrojenie LLM do monitorowania własnej wiarygodności \cite{predmetacog2026} poprawia kalibrację, ale samo w sobie nie czyni aktualizacji online, więc $\Gamma$ pozostaje zerowe przy wdrożeniu.

\textbf{Titans} \cite{behrouz2024}, moduł neuronowej pamięci długoterminowej aktualizowany w czasie inferencji przez sygnał ,,zaskoczenia'' (gradient straty pamięci asocjacyjnej względem wejścia) napędzający aktualizację online $S_t=\arg\min_S\frac12\|Sk_t-v_t\|^2+\frac1{2\eta}\|S-S_{t-1}\|^2$, jest konkretną, działającą instancją $e_\ell(t)>0$ napędzającego $g_\ell(t)>0$. Jednak Titans przewiduje \emph{treść kontekstu} (klucze/wartości z wejścia), nie własny stan wewnętrzny systemu --- realna niejasność, którą rozstrzygamy poniżej.

\subsection*{Rozstrzygnięcie self- vs. world-modeling: dwa niezależne warunki}
(a) \emph{Cel}: czy przewidywana wielkość to wewnętrzny stan reprezentacyjny ($h_\ell$, aktywacje ukryte) czy zewnętrznie obserwowalne zachowanie (token) --- nawet jeśli ten token jest autorstwa samego systemu? (b) \emph{Moment}: czy aktualizacja zachodzi online przy wdrożeniu, czy tylko podczas treningu (po którym wagi zamrażają się)?

\begin{center}
\small
\begin{tabular}{@{}>{\raggedright}p{2.6cm}>{\raggedright}p{4.5cm}>{\raggedright\arraybackslash}p{4.5cm}@{}}
\toprule
 & \textbf{tylko trening} & \textbf{online przy wdrożeniu} \\
\midrule
cel = zachowanie zewnętrzne & zwykły trening predykcyjny & \textbf{Titans}: konstytutywne, ale \emph{world}-modeling \\[4pt]
cel = stan wewnętrzny & self-modelling jako zadanie pomocnicze \cite{premakumar2026} & \textbf{Temporal Self-Model} \cite{tsm2026}: pełna zgodność \\
\bottomrule
\end{tabular}
\end{center}
Tylko prawa-dolna komórka spełnia naszą pełną definicję konstytutywnego self-modelingu. Titans, mimo rozgłosu, ląduje gdzie indziej. \cite{premakumar2026} pokazują sieć trenowaną z pomocniczym celem samo-predykcji stanu ukrytego, która staje się prostsza i bardziej zregularyzowana --- ale aktualizacja jest zamrożona po treningu, więc deskryptywna przy wdrożeniu. \cite{tsm2026} opisuje moduł jawnie przewidujący własne przyszłe stany ukryte agenta, z opóźnionym porównaniem jako błędem samo-predykcji --- najbliższe naszej definicji, choć czy jego aktualizacja jest rzeczywiście online (nie tylko trenowana standardowym backpropem) wymaga dalszego potwierdzenia.

Przypadek pośredni to \emph{in-context learning jako niejawny spadek gradientu} \cite{oswald2023}: pojedynczy przebieg w przód przez uwagę może zaimplementować krok strukturalnie równoważny spadkowi gradientu po wewnętrznym celu, ale wywołana zmiana nie przetrwa poza oknem kontekstu (brak trwałej zmiany $\theta$) --- ani w pełni deskryptywny, ani w pełni konstytutywny w naszym sensie.

Osobno: \textbf{warstwa homeostatyczna} jest potwierdzonym, realnym brakiem w mainstreamowych architekturach transformerowych; ramy regulacyjne inspirowane interocepcją dla AI \cite{interoceptive2026} istnieją tylko jako nurt marginalny. \textbf{Prawdziwa rekurencja} (trwały stan wewnętrzny ewoluujący niezależnie od nowych tokenów) jest obecna w modelach przestrzeni stanów (Mamba) i hybrydach typu Titans, nieobecna w czystej uwadze (KV-cache to pamięć kontekstu, nie pętla dynamiczna).

\section{Hipoteza: dwa niezależne mechanizmy znieczulenia w $(\Sigma,\Gamma,A)$}\label{sec:anesthesia}

Przy stałym poziomie integracji samoreferencyjność można zredukować na dwa jakościowo różne sposoby: (1) $\Gamma$ spada bezpośrednio, w tym samym miejscu, albo (2) przesuwa się w obszar o niższej integracji (spadek $A$, niekoniecznie $\Gamma$). Literatura anestezjologiczna dostarcza dwóch osobnych, pasujących do tego rozróżnienia przypadków.

\textbf{Mechanizm 1 ($\Gamma$ spada w miejscu).} Propofol i sewofluran uderzają w tzw. \emph{posterior hot zone} \cite{hotzone2021} --- region, który zwolennicy IIT wprost wskazują jako główny substrat świadomości (w przeciwieństwie do GNW, stawiającego na sieci czołowo-ciemieniowe) --- i jednocześnie zaburzają sprzężenie zwrotne czołowo-ciemieniowe (front$\to$back) \cite{lee2013disruption}, związane z przetwarzaniem wyższego rzędu. Integracja i samo-korekta zapadają się razem, w tym samym miejscu. Behawioralnie: pełna utrata przytomności, brak jakiegokolwiek raportu --- spójne z $\Sigma$ i $\Gamma$ spadającymi razem.

\textbf{Mechanizm 2 ($\Gamma$ zostaje, $A$ spada).} Ketamina w dawkach subanestetycznych daje dysocjację --- reaktywność i często żywe, bogate doświadczenie mogą się utrzymać, ale z zaburzonym poczuciem siebie. Rozregulowuje osobno sieć wyrazistości (,,minimalne'', cielesne self) i osobno DMN (,,biograficzne'', narracyjne self) \cite{ketaminemodel2022} --- dwa piętra samo-modelowania, dość dobrze odpowiadające warstwom $\ell$ z tego modelu. Ogólna bogatość doświadczenia nie zapada się tak jak przy propofolu (stąd żywe, ,,psychodeliczne'' jakościowo relacje z ketaminy) --- samo-korekta nie znika, tylko odłącza się od miejsca, gdzie akurat jest integracja: spadek $A$, niekoniecznie $\Sigma$ czy $\Gamma$.

\textbf{Status:} wiarygodna, umocowana w anatomii hipoteza, nie potwierdzony wynik --- nikt nie policzył $(\Sigma,\Gamma,A)$ wprost na danych EEG/fMRI z obu anestetyków. Konkretna, testowalna przepowiednia: profil $(\Sigma,\Gamma,A)$ powinien wyraźnie różnić propofol/sewofluran ($\Sigma\downarrow,\Gamma\downarrow$ razem) od ketaminy ($\Sigma,\Gamma$ względnie zachowane, $A\downarrow$), mimo że oba ,,wyłączają'' jakąś formę normalnego funkcjonowania.

\section{Kanał odczytu: rozdzielenie stanu wewnętrznego od obserwowalności}\label{sec:readout}

Hipoteza o znieczuleniu powyżej (\S\ref{sec:anesthesia}) odsłania oś, której formalizm wcześniej nie miał: $(\Sigma,\Gamma,A)$ opisuje stan \emph{wewnętrzny}, ale nic nie mówi, czy ten stan da się odczytać z zachowania. To dokładnie problem kliniczny, który zmotywował powstanie Perturbational Complexity Index --- miary jawnie zaprojektowanej jako ,,niezależna od przetwarzania sensorycznego i zachowania'' \cite{casali2013}.

\subsection*{Definicja}
Niech $M(t)$ będzie osobnym substratem eferentnym (jednostki ruchowe/komunikacyjne), odrębnym od warstw $\ell$. Definiujemy kanał odczytu z warstwy $\ell$ do $M$, tą samą miarą odległości $D$ co bazowe $\Phi$:
\[
R_\ell(t) = D\big(p(m'\mid s_\ell,m)\,\|\,p(m'\mid m)\big),
\]
tj. o ile znajomość stanu warstwy $\ell$ zmienia przewidywalność przyszłego stanu efektorów, ponad to, co daje sama znajomość aktualnego stanu efektorów. Strukturalnie identyczne z $\Phi_{\ell,\ell-1}$ (\S3.1), tylko skierowane na kanał warstwa$\to$wyjście. Szczególnie istotne jest $R_L(t)$ --- kanał z najwyższej, samo-modelującej warstwy do efektorów.

\begin{proposition}[Własności kanału odczytu]
$R_\ell(t)\ge0$ (dywergencja), oraz $R_\ell(t)=0$ wtedy i tylko wtedy, gdy $s_\ell(t)\perp m'\mid m(t)$ --- odczyt znika dokładnie wtedy, gdy warstwa jest warunkowo nieistotna dla przyszłego stanu efektorów przy danym aktualnym, tj. nie ma w ogóle kanału do wyjścia. Jeśli $D$ odczytać jako dywergencję KL (więc $R_\ell=I(s_\ell;m'\mid m)$), zachodzi nierówność przetwarzania danych: jeśli każda ścieżka przyczynowa z warstwy $\ell$ do $M$ przechodzi przez pewną warstwę $\ell'$, to $R_\ell\le R_{\ell'}$.
\end{proposition}

\subsection*{Obserwowalne vs.\ rzeczywiste}
To, co widzi klinicysta ($\Sigma_{\text{obs}}$ --- ,,widoczna'' świadomość, wnioskowana z zachowania) jest wąskim gardłem przez $R_L$, niezależnie od faktycznego $\Sigma$:
\[
\Sigma_{\text{obs}}(t) \le \beta\big(R_L(t)\big), \qquad \beta:[0,\infty)\to[0,\infty)\ \text{niemalejąca},\ \beta(0)=0,
\]
poziom widoczny zdominowany przez niemalejącą funkcję odczytu, znikającą razem z nim. Kształt $\beta$ zostaje otwarty (zależy od punktacji zachowania); nośna konsekwencja to punkt końcowy: $R_L(t)=0 \Rightarrow \Sigma_{\text{obs}}(t)=0$ \emph{niezależnie od tego, jak wysokie jest rzeczywiste $\Sigma(t)$} --- formalizacja ryzyka mylenia $\Sigma_{\text{obs}}$ z $\Sigma$. (Wcześniejsza wersja robocza pisała $\Sigma_{\text{obs}}\le\kappa\cdot R_L$ ,,dla pewnego $\kappa>0$'', co jest jałowe.)

\subsection*{Mapowanie kliniczne}
\textbf{Zespół zamknięcia:} $\Sigma,\Gamma,A$ normalne; $R_L\approx0$ z powodu strukturalnego uszkodzenia dróg eferentnych (np. udar brzusznej części mostu) $\to\Sigma_{\text{obs}}\approx0$ mimo wysokiego $\Sigma$. \textbf{Sen REM:} $\Sigma,\Gamma,A$ względnie zachowane (PCI$\approx$czuwanie); $R_L$ aktywnie stłumione przez znany, odwracalny mechanizm fizjologiczny --- atonia REM, hamowanie neuronów ruchowych rdzenia przez jądro podsinawe --- potwierdzając, że $R_\ell$ ma realny substrat neuronalny, nie tylko teoretyczną wygodę. \textbf{Cognitive motor dissociation} (część pacjentów w stanie wegetatywnym, paradygmaty Owena): $R_L$ częściowe --- wystarcza na modulację odpowiedzi w neuroobrazowaniu na polecenie, nie wystarcza do mięśni. \textbf{Prawdziwy stan wegetatywny (bez ukrytej świadomości):} i $\Sigma$, i $R_L$ niskie --- różnica, którą PCI miało rozstrzygnąć.

\section{Wrażliwość kontekstowa bramki konstytutywnej: context blindness i ,,zbyt intensywny świat''}

Źródło omawiające zasadę wolnej energii Fristona w kontekście autyzmu \cite{vandecruys2014} formułuje to ostro: przetwarzanie kontekstu wymaga, żeby mózg nadał swoim oczekiwaniom odpowiednią wagę; jeśli oczekiwania dostają zbyt małą wagę, albo błędy predykcji pozostają zbyt sztywno duże, integrowanie kontekstu zawodzi. Van de Cruys i in. proponują \emph{aberrant precision} jako mechanizm autyzmu --- błędy predykcji dostają zbyt sztywno wysoką wagę, niezależnie od kontekstu. \cite{qela2025} mapują pokrewne odchylenia mechanizmu predykcyjnego w autyzmie, schizofrenii i depresji.

\subsection*{Domykająca luka w modelu}
,,Waga nadawana oczekiwaniom'' to w żargonie predictive processing \emph{precyzja} --- strukturalnie dokładnie nasze $g_\ell(t)$. Formalizm nie miał dotąd miejsca na to, żeby \emph{kontekst} normalnie tę wagę modulował. Rozszerzenie:
\[
g_\ell(t) = h\big(e_\ell(t), C(t)\big), \qquad \chi_\ell(t) = \left|\frac{\partial g_\ell}{\partial C}\right| \ge 0,
\]
gdzie $C(t)$ to sygnał kontekstowy, a $\chi_\ell(t)$ --- nowy parametr --- mierzy, jak silnie kontekst moduluje wagę błędu. $\chi_\ell(t)\ge0$ trywialnie.

\begin{proposition}[Ograniczona całkowita modulacja]
Ponieważ $g_\ell\in[0,1]$, wzdłuż dowolnej ścieżki kontekstu $C_0\to C_1$, na której $g_\ell$ jest monotoniczne, $\int_{C_0}^{C_1}\chi_\ell\,dC = |g_\ell(C_1)-g_\ell(C_0)|\le 1$. Wrażliwość kontekstowa to skończony budżet: wysokie $\chi_\ell$ na jednym odcinku wymusza niskie $\chi_\ell$ (względną context blindness) gdzie indziej, gdy $g_\ell$ nasyci się przy $0$ lub $1$.
\end{proposition}

\subsection*{Dwa reżimy}
\textbf{Normalna, kontekstowo czuła precyzja} ($\chi_\ell>0$): znane, kontekstowo nieistotne sygnały dostają automatycznie niższą wagę --- $g_\ell$ spada niezależnie od surowej wielkości $e_\ell$. \textbf{Context blindness} ($\chi_\ell\approx0$): $g_\ell(t)\approx\hat g_\ell(e_\ell(t))$, sztywna funkcja samego błędu; kontekst nic nie zmienia. Efekt: sygnały, które kontekst normalnie by wyciszył, dalej dostają pełną wagę --- formalny odpowiednik ,,zbyt intensywnego świata'' (Intense World Theory, Markram \& Markram 2010): nic nie jest domyślnie tłumione jako ,,już przewidziane w tym kontekście''.

\subsection*{Konsekwencja dla $\Gamma$}
W typowym przypadku powtarzana, kontekstowo nieistotna stymulacja obniża i $e_\ell$ (przez uczenie predykcyjne --- repetition suppression), i $g_\ell$ (przez aktywne tłumienie kontekstowe, $\chi_\ell>0$) --- podwójny spadek wkładu do $\Gamma$. Przy context blindness działa tylko pierwszy mechanizm (jeśli w ogóle) --- $\Gamma(t)$ pozostaje podwyższone/zmienne nawet w środowiskach o niskiej nowości. Formalizowalna, przynajmniej jakościowo testowalna różnica dynamiki $\Gamma$ między typowym a ,,aberrant precision'' przypadkiem --- nikt jej jednak nie zmierzył wprost w tych terminach.

\section{Dalsze rozważania (nieformalne --- bez dowodów, niższy status niż reszta pracy)}

\subsection*{Skala: od pojedynczej komórki po cały wszechświat}

\textbf{Dolna granica.} Czyste IIT jest liberalne: ,,Minimal physicalism as a scale-free substrate for cognition and consciousness'' (2021) argumentuje, że sieci regulacji genów i transdukcji sygnału na poziomie pojedynczej komórki mają dodatnie $\Phi$. \textbf{Nasz model jest bardziej restrykcyjny} --- dodatnie $\Phi$ nie wystarcza, potrzeba $\Gamma>0$, realnej, online działającej pętli samo-predykcji. Pojedyncza komórka ma homeostazę, nie self-model w sensie \S3. \emph{C. elegans} (302 neurony, konektom w większości feedforward, niewielka rekurencja) to dobry przypadek testowy: prawdopodobnie niskie/zerowe $\Gamma$ mimo niezerowego $\Phi$.

\textbf{Górna granica.} ,,Upper bounds for integrated information'' (2024) pokazuje, że $\Phi$ może rosnąć hiperwykładniczo z liczbą jednostek --- matematycznie brak sufitu. Ale: (1) postulat wykluczenia nie pozwala na sumowanie się w górę --- większy system wygrywa tylko jeśli ma \emph{wyższe} $\Phi$, nie więcej jednostek; (2) systemy wielkoskalowe (internet, biosfera, społeczeństwo) mają połączenia rzadkie i wolne względem neuronów, więc empirycznie nie przebijają mózgu.

\textbf{Cały wszechświat --- ,,podwójne nie''.} To dosłownie \emph{problem kombinacji} panpsychizmu, doprowadzony do granicy. Dwie niezależne linie rozumowania zbiegają się w tym samym wniosku: (1) \emph{strukturalnie}, odpowiedź Kocha/Tononiego na zarzut Searle'a (,,świadomość rozmazana jak dżem'') --- ,,moja świadomość, twoja świadomość, ale nic pomiędzy'' --- wszechświat to mozaika lokalnych maksimów (,,ontologiczny pył''), nigdy nadrzędna całość, bo agregat (np. mózg + reszta kosmosu) ma niższe $\Phi$ niż jego gęściej połączone części; (2) \emph{fizycznie}, regiony poza swoimi wzajemnymi horyzontami kosmologicznymi nigdy nie były w kontakcie przyczynowym --- nie ma nawet czego wstawić do macierzy przejść, którą wymaga $\Phi$. Pod $\Sigma$ dochodzi trzeci powód: brak jakiegokolwiek jednolitego $\Gamma$ na skalę kosmologiczną. Ciekawostka na marginesie, bez rygoru: skoro integracja wymaga gęstej, szybkiej struktury, a wszechświat się rozrzedza w stronę śmierci cieplnej, możliwe że jesteśmy bliżej szczytu możliwej gęstości integracji w historii kosmosu niż w jego dalekiej przyszłości.

\textbf{Zastrzeżenie nadrzędne:} samo IIT, na którym to wszystko się opiera, jest naukowo kontestowane (2023: list otwarty 124 badaczy, w tym Dennett, nazywający IIT pseudonauką; COGITATE/Templeton --- wynik niejednoznaczny wobec GNW).

\subsection*{Zespół obcej ręki jako ilustracja netto fragmentacji (\S5)}

Najczystsza znaleziona ilustracja $\Sigma<0$. Przy uszkodzeniu spoidła wielkiego (corpus callosum) --- obniżonym $\Phi_{\ell,\ell-1}$ między półkulami --- dwa kandydackie kompleksy motoryczno-intencjonalne nie tylko przestają współpracować, ale \textbf{aktywnie się sobie przeciwstawiają}: ,,diagonistic dyspraxia'' --- jedna ręka dosłownie cofa to, co zrobiła druga. To integracja poniżej zera, nie tylko jej brak --- dokładnie mechanizm z \S5, i ta sama linia badań (split-brain, corpus callosum) co interpreter Gazzanigi motywujący \S3 na samym początku.

\textbf{Słabsze, częściowe dopasowanie: dysocjacyjne zaburzenie tożsamości (DID).} Dobrze udokumentowane: odrębne, skompartmentalizowane sygnatury neuronalne dla różnych stanów tożsamości, przesunięcia łączności przy przełączeniach --- pasuje do $\ell^*(t)$ zmieniającego się w czasie (scenariusz dysocjacji z \S6). Słabiej udokumentowane: czy któryś stan \textbf{aktywnie degraduje} integrację innego (prawdziwe $\Sigma<0$), czy tylko go zastępuje (bariery amnezji to bardziej ,,brak dostępu'' niż aktywna interferencja w sensie $w_-$). Status DID jako kategorii diagnostycznej jest też przedmiotem realnej debaty naukowej (model traumatyczny vs. socjokognitywny) --- nierozstrzygniętej.

\section{Dyskusja i otwarte pytania}



To jest spójna wewnętrznie synteza koncepcyjna, nie empirycznie zweryfikowana teoria. W sprawie mechanizmu wykluczenia (\S\ref{sec:exclusion}) wybór jest przesądzony: kanoniczna jest miękka postać $w_+-w_-$, a osiągalne $\Sigma<0$ (netto fragmentacja) zostaje jako zamierzona cecha, nie usterka; binarna alternatywa z $\Sigma\ge0$ pozostaje wyłącznie jako udokumentowany, odrzucony wariant. Otwarte punkty: systematyczny (nie tylko oparty na wyszukiwarce) przegląd literatury; oraz empiryczne zastosowanie na architekturach poza przykładem zabawkowym, zwłaszcza żeby sprawdzić, czy trajektorie $(\Sigma,\Gamma,A)$ śledzą znane markery integracji (np. Perturbational Complexity Index pod znieczuleniem).

% Cytowania zweryfikowane ze źródłami pierwotnymi, 2026-08-28 (dwie tury).
% Kilka numerów tomów/artykułów z drugiej tury wymaga jeszcze potwierdzenia
% u wydawcy --- patrz notes/literature.md.
\begin{thebibliography}{99}
\bibitem{aaronson2014} Aaronson, S. (2014). Why I am not an integrated information theorist (or, the unconscious expander). \emph{Shtetl-Optimized} (blog), 21 maja 2014. \url{https://scottaaronson.blog/?p=1799}
\bibitem{albantakis2023} Albantakis, L., Barbosa, L., Findlay, G., Grasso, M., Haun, A. M., Marshall, W., Mayner, W. G. P., Zaeemzadeh, A., Boly, M., Juel, B. E., Sasai, S., Fujii, K., David, I., Hendren, J., Lang, J. P., \& Tononi, G. (2023). Integrated information theory (IIT) 4.0: Formulating the properties of phenomenal existence in physical terms. \emph{PLOS Computational Biology}, 19(10), e1011465.
\bibitem{riddle2024} Riddle, J., \& Schooler, J. W. (2024). Hierarchical consciousness: the Nested Observer Windows model. \emph{Neuroscience of Consciousness}, 2024(1), niae010.
\bibitem{usk2026} Tallam, K. (2026). Consciousness as uncommon self-knowledge: A synergistic information framework. \emph{arXiv:2605.13884}.
\bibitem{friston2010} Friston, K. (2010). The free-energy principle: a unified brain theory? \emph{Nature Reviews Neuroscience}, 11(2), 127--138.
\bibitem{deane2021} Deane, G. (2021). Consciousness in active inference: Deep self-models, other minds, and the challenge of psychedelic-induced ego-dissolution. \emph{Neuroscience of Consciousness}, 2021(2), niab024.
\bibitem{rosenthal2005} Rosenthal, D. M. (2005). \emph{Consciousness and Mind}. Oxford University Press.
\bibitem{behrouz2024} Behrouz, A., Zhong, P., \& Mirrokni, V. (2024). Titans: Learning to memorize at test time. \emph{arXiv:2501.00663}.
\bibitem{oswald2023} von Oswald, J., Niklasson, E., Randazzo, E., Sacramento, J., Mordvintsev, A., Zhmoginov, A., \& Vladymyrov, M. (2023). Transformers learn in-context by gradient descent. \emph{Proceedings of the 40th International Conference on Machine Learning (ICML)}, PMLR 202.
\bibitem{premakumar2026} Premakumar, V. N., Vaiana, M., Pop, F., Rosenblatt, J., Schwerz de Lucena, D., Ziman, K., \& Graziano, M. S. A. (2026). Unexpected benefits of self-modelling in neural systems. \emph{Philosophical Transactions of the Royal Society A}, 384(2320), 20240531. Preprint: arXiv:2407.10188 (2024).
\bibitem{tsm2026} Xie, Y. (2026). Self-monitoring benefits from structural integration: Lessons from metacognition in continuous-time multi-timescale agents. \emph{arXiv:2604.11914}.
\bibitem{interoceptive2026} Candia-Rivera, D. (2026). Interoceptive machine framework: Toward interoception-inspired regulatory architectures in artificial intelligence. \emph{arXiv:2604.24527}.
\bibitem{casali2013} Casali, A. G., Gosseries, O., Rosanova, M., Boly, M., Sarasso, S., Casali, K. R., Casarotto, S., Bruno, M.-A., Laureys, S., Tononi, G., \& Massimini, M. (2013). A theoretically based index of consciousness independent of sensory processing and behavior. \emph{Science Translational Medicine}, 5(198), 198ra105.
\bibitem{cerullo2015} Cerullo, M. A. (2015). The problem with Phi: A critique of integrated information theory. \emph{PLOS Computational Biology}, 11(9), e1004286.
\bibitem{hotzone2021} Ihalainen, R., Gosseries, O., Van de Steen, F., Raimondo, F., Panda, R., Bonhomme, V., Marinazzo, D., Bowman, H., Laureys, S., \& Chennu, S. (2021). How hot is the hot zone? Computational modelling clarifies the role of parietal and frontoparietal connectivity during anaesthetic-induced loss of consciousness. \emph{NeuroImage}, 231, 117841.
\bibitem{lee2013disruption} Lee, U., Ku, S., Noh, G., Baek, S., Choi, B., \& Mashour, G. A. (2013). Disruption of frontal--parietal communication by ketamine, propofol, and sevoflurane. \emph{Anesthesiology}, 118(6), 1264--1275.
\bibitem{ketaminemodel2022} Marguilho, M., Figueiredo, I., \& Castro-Rodrigues, P. (2023). A unified model of ketamine's dissociative and psychedelic properties. \emph{Journal of Psychopharmacology}, 37(1), 14--32. (Opublikowano online w grudniu 2022.)
\bibitem{vandecruys2014} Van de Cruys, S., Evers, K., Van der Hallen, R., Van Eylen, L., Boets, B., de-Wit, L., \& Wagemans, J. (2014). Precise minds in uncertain worlds: Predictive coding in autism. \emph{Psychological Review}, 121(4), 649--675.
\bibitem{qela2025} Qela, B., Damiani, S., De Santis, S., i in. (2025). Predictive coding in neuropsychiatric disorders: A systematic transdiagnostic review. \emph{Neuroscience \& Biobehavioral Reviews}, 169, 106020.
\bibitem{bayne2016} Bayne, T., Hohwy, J., \& Owen, A. M. (2016). Are there levels of consciousness? \emph{Trends in Cognitive Sciences}, 20(6), 405--413.
\bibitem{safron2020} Safron, A. (2020). An Integrated World Modeling Theory (IWMT) of consciousness: Combining integrated information and global neuronal workspace theories with the free energy principle and active inference framework. \emph{Frontiers in Artificial Intelligence}, 3, 30.
\bibitem{hanson2023} Hanson, J. R., \& Walker, S. I. (2023). On the non-uniqueness problem in integrated information theory. \emph{Neuroscience of Consciousness}, 2023(1), niad014.
\bibitem{barrett2019} Barrett, A. B., \& Mediano, P. A. M. (2019). The $\Phi$ measure of integrated information is not well-defined for general physical systems. \emph{arXiv:1902.04321}.
\bibitem{mediano2019} Mediano, P. A. M., Seth, A. K., \& Barrett, A. B. (2019). Measuring integrated information: Comparison of candidate measures in theory and simulation. \emph{Entropy}, 21(1), 17.
\bibitem{phiid2019} Mediano, P. A. M., Rosas, F. E., Carhart-Harris, R. L., Seth, A. K., \& Barrett, A. B. (2019). Beyond integrated information: A taxonomy of information dynamics phenomena. \emph{arXiv:1909.02297}.
\bibitem{oizumi2016} Oizumi, M., Amari, S., Yanagawa, T., Fujii, N., \& Tsuchiya, N. (2016). Measuring integrated information from the decoding perspective. \emph{PLOS Computational Biology}, 12(1), e1004654.
\bibitem{intrepid2026} Corcoran, A. W., Haun, A. M., Dorman, R., Tononi, G., Friston, K. J., Pennartz, C. M. A., \& the TWCF:INTREPID Consortium. (2026). Integrated information and predictive processing theories of consciousness: An adversarial collaborative review. \emph{Neuroscience \& Biobehavioral Reviews}, 187, 106742. Preprint: arXiv:2509.00555.
\bibitem{predmetacog2026} Luo, W., \& Jho, H. (2026). Predictive metacognition: A neuro-computational framework for self-monitoring in large language models. \emph{Scientific Reports}. \url{https://doi.org/10.1038/s41598-026-54840-2}
\end{thebibliography}

\end{document}
